% =============================================================================
%  Datei-Hinweise & Konfiguration: siehe config.md
%  -> Abschnitt "anhang/ki-prompts.tex"
% =============================================================================

\chapter{Anhang: Verwendete KI-Prompts}
\label{cha:anhang-prompts}

Im Verlauf der Erstellung dieser Arbeit wurden vier generative KI-Werkzeuge herangezogen. Anthropic Claude (Opus~5) diente der Erstellung des Glossars, des Formelzeichenverzeichnisses und des Anhangs, dem Entwurf der Marktdatenerfassung und ihrer Persistenz, dem Abgleich der Schnittstellendokumentation der Bank sowie der Anpassung des Extraktionsskripts für die Code- und Datenquellen der Cloud-Umgebung, der Erzeugung der Testskripte zur Ergebniskontrolle und der Testsuite für die Betriebsskripte, die Ausführungspfade und die Mutationsprobe. Anthropic Claude (Fable~5) unterstützte bei den Python- und Shell-Skripten des Agenten, bei der Cloud-Run-Anwendung des öffentlichen Dashboards, beim Totmann-Schalter, bei der statistischen Prüfinstanz samt Prognosemodul, bei der Gegenprüfung der Testskripte sowie bei der Erzeugung der beiden als KI-generiert gekennzeichneten Abbildungen. Google Gemini (3.5~Flash) wurde für die Rechtschreib- und Interpunktionsprüfung, die Quellenrecherche, den Abgleich der automatischen Tests und die Erstellung des Abkürzungsverzeichnisses eingesetzt. Moonshot AI Kimi (K3) diente einem allgemeinen Codereview der beigelegten Quellcodebeilage und dem Erstentwurf des Änderungsverzeichnisses der Beilage.

Gemäß FOM-Leitfaden Kap.~6.2.4 sind die genutzten Prompts in Kurzform — ohne die jeweils generierten Antworten — zu dokumentieren. Sämtliche Antworten wurden inhaltlich geprüft, eigenständig verifiziert und nicht wörtlich in die Arbeit übernommen.

% \kiwerkzeug{Tool}{Version/Stand}{Zugriffs-URL} -- einheitliche Darstellung
% eines KI-Werkzeugs; das zugehoerige \nocite{<bibkey>} steht in der
% jeweiligen Tool-Datei (Eintrag im KI-Verzeichnis, Kap. 6.2.4).
\newcommand{\kiwerkzeug}[3]{%
    \paragraph*{Verwendetes Tool}
    \begin{description}[leftmargin=2.5cm,style=nextline,labelwidth=2cm]
        \item[Tool] #1
        \item[Version] #2
        \item[Zugriff] \url{#3}
    \end{description}%
}

% --- Eingesetzte KI-Werkzeuge (je Werkzeug eine modulare Datei) -------------
% REIHENFOLGE NICHT AENDERN: Sie entspricht der Sortierung des KI-Verzeichnisses
% (biblatex sortiert nach Urheber: Anthropic, Anthropic, Google DeepMind,
% Moonshot AI). Die \note-Felder der Bib-Eintraege verweisen mit Anhang III.I
% bis III.IV auf genau diese Abschnittsnummern.
%  =============================================================================
%  Datei-Hinweise & Konfiguration: siehe config.md
%  -> Abschnitt "anhang/ki-prompts/tool-claude.tex"
%  =============================================================================

\section{Anthropic Claude (Opus~5)}
\label{sec:prompts-claude}

\kiwerkzeug{Anthropic Claude, Modell Opus 5}{Modellbezeichner
\texttt{claude-opus-5}, veröffentlicht am 24.07.2026}{https://claude.ai}
% Eintrag im KI-Verzeichnis (Leitfaden Kap. 6.2.4) erzwingen mittels \nocite.
\nocite{claudeOpus2026}

Eingesetzt für die Verzeichnisse des Vorspanns und den Aufbau des Anhangs
sowie für zwei Werkzeuge der Projektumgebung: die Anpassung des
Extraktionsskripts, das Quellcode- und Datenbestände der Google-Cloud-Umgebung
für das elektronische Archiv aufbereitet, und die Erzeugung der Testskripte
zur Ergebniskontrolle.

\subsection{Erstellung des Glossars}
\label{sec:prompts-claude-glossar}
\begin{enumerate}
    \item \enquote{Formuliere zu den folgenden Fachbegriffen jeweils eine
            prägnante, allgemeinverständliche Glossardefinition von höchstens
            zwei Sätzen im sachlichen wissenschaftlichen Stil: [\dots]}
            \par\hfill\textit{Kontext: Glossar (Vorspann)}

    \item \enquote{Prüfe die folgenden Glossareinträge auf fachliche
            Korrektheit, Widerspruchsfreiheit und einheitliche Formulierung:
            [\dots]}
            \par\hfill\textit{Kontext: Glossar (Vorspann)}
\end{enumerate}

\subsection{Erstellung des Formelzeichenverzeichnisses}
\label{sec:prompts-claude-formelzeichen}
\begin{enumerate}
    \item \enquote{Extrahiere aus den folgenden nummerierten Gleichungen
            sämtliche Formelzeichen und formuliere je Zeichen eine
            Kurzerklärung, die Bedeutung und Wertebereich beziehungsweise
            Einheit benennt: [\dots]}
            \par\hfill\textit{Kontext: Formelzeichenverzeichnis (Vorspann)}

    \item \enquote{Prüfe, ob jedes in den Gleichungen verwendete Formelzeichen
            im Verzeichnis geführt wird und ob umgekehrt jeder Eintrag des
            Verzeichnisses im Text tatsächlich vorkommt; nenne Abweichungen in
            beide Richtungen: [\dots]}
            \par\hfill\textit{Kontext: Formelzeichenverzeichnis (Vorspann)}
\end{enumerate}

\subsection{Aufbau des Anhangs}
\label{sec:prompts-claude-anhang}
\begin{enumerate}
    \item \enquote{Schlage eine Gliederung für einen dreiteiligen Anhang vor,
            der ein elektronisches Quellcodearchiv, die Belege einer
            Evaluation und die verwendeten KI-Prompts umfasst; benenne je
            Teil, welche Angaben erforderlich sind, damit ein Dritter die
            Ergebnisse nachvollziehen kann: [\dots]}
            \par\hfill\textit{Kontext: Anhang~\ref{cha:anhang-archiv} bis
            Anhang~\ref{cha:anhang-prompts}}

    \item \enquote{Formuliere zu den nachfolgenden Rohprotokollen jeweils
            einen einleitenden Absatz, der Herkunft, Erzeugungszeitpunkt und
            Aussagegehalt des Belegs benennt, ohne die Zahlen selbst zu
            wiederholen: [\dots]}
            \par\hfill\textit{Kontext:
            Anhang~\ref{sec:anhang-evaluationsbelege}}
\end{enumerate}

\subsection{Anpassung des Extraktionsskripts für Code- und Datenquellen}
\label{sec:prompts-claude-extraktion}

Das Skript zur Aufbereitung des Quellcodearchivs lag aus vorangegangenen
Projekten bereits vor; Gegenstand der Prompts war ausschließlich seine
Anpassung an die Zielarchitektur dieser Arbeit (\cref{fig:zielarchitektur})
mit ihren zwei virtuellen Maschinen, Journaldateien und zeitgesteuerten
Aufträgen.

\begin{enumerate}
    \item \enquote{Erweitere das vorhandene Extraktionsskript so, dass es
            neben Python- und Shell-Quelltexten auch die
            \texttt{systemd}-Einheiten und Cron-Definitionen beider
            virtueller Maschinen erfasst und je Datei Herkunftsmaschine sowie
            Ausführungsintervall in einem Manifest festhält: [\dots]}
            \par\hfill\textit{Kontext:
            Anhang~\ref{sec:anhang-archiv-quellcode}}

    \item \enquote{Ergänze das Skript um eine Redaktionsstufe, die
            Zugangsdaten sowie Konto- und Ordernummern aus den Journaldateien
            entfernt, bevor diese in das Archiv aufgenommen werden, und die
            jede Ersetzung protokolliert: [\dots]}
            \par\hfill\textit{Kontext: Anhang~\ref{sec:anhang-archiv-daten}}

    \item \enquote{Stelle sicher, dass die erzeugte Archivdatei bei
            unverändertem Inhalt reproduzierbar dieselbe Prüfsumme liefert,
            und erläutere, welche Zeitstempel dafür festzusetzen sind:
            [\dots]}
            \par\hfill\textit{Kontext:
            Anhang~\ref{sec:anhang-archiv-identitaet}}
\end{enumerate}

\subsection{Erzeugung der Testskripte zur Ergebniskontrolle}
\label{sec:prompts-claude-tests}
\begin{enumerate}
    \item \enquote{Entwirf zu der nachfolgenden Auswertungsfunktion eine
            Testsuite mit gemockten Schnittstellen, die jede
            sicherheitsrelevante Verzweigung abdeckt -- insbesondere die
            Pfade, die zu einer Ablehnung oder zu einem Handelsstopp führen:
            [\dots]}
            \par\hfill\textit{Kontext: \cref{sec:umsetzung-implementierung}}

    \item \enquote{Formuliere Testfälle, die prüfen, ob die im Protokoll
            ausgewiesenen Kennzahlen aus denselben Eingangsdaten
            reproduzierbar hervorgehen, und benenne die Toleranz, innerhalb
            derer Gleitkommaabweichungen zulässig sind: [\dots]}
            \par\hfill\textit{Kontext: \cref{sec:umsetzung-evaluation}}
\end{enumerate}

%  =============================================================================
%  Ende
%  =============================================================================

%  =============================================================================
%  Datei-Hinweise & Konfiguration: siehe config.md
%  -> Abschnitt "anhang/ki-prompts/tool-fable.tex"
%  =============================================================================

\section{Anthropic Claude (Fable~5)}
\label{sec:prompts-fable}

\kiwerkzeug{Anthropic Claude, Modell Fable 5}{Modellbezeichner
\texttt{claude-fable-5}, veröffentlicht am 09.06.2026}{https://claude.ai}
% Eintrag im KI-Verzeichnis (Leitfaden Kap. 6.2.4) erzwingen mittels \nocite.
\nocite{claudeFable2026}

Eingesetzt zur Unterstützung bei der Implementierung: Python- und
Shell-Skripte des Agenten, die Cloud-Run-Anwendung des öffentlichen
Dashboards, die Gegenprüfung der zuvor erzeugten Testskripte sowie die beiden
als KI-generiert gekennzeichneten Abbildungen dieser Arbeit
(\cref{fig:zielarchitektur} und \cref{fig:projektstruktur}).

\subsection{Unterstützung bei Python- und Shell-Skripten}
\label{sec:prompts-fable-skripte}
\begin{enumerate}
    \item \enquote{Überarbeite die folgende Funktion so, dass ein
            Verbindungsabbruch des Datenstroms nicht zum Prozessabbruch führt,
            sondern erkannt, protokolliert und über einen Wächterprozess
            neu verbunden wird; erhalte dabei die bestehende Signatur:
            [\dots]}
            \par\hfill\textit{Kontext: \cref{sec:umsetzung-architektur}}

    \item \enquote{Schreibe ein Shell-Skript, das den Entscheidungszyklus
            startet, eine Prozesssperre je Konto setzt, überlappende Läufe
            verhindert und im Fehlerfall mit eindeutigem Rückgabewert endet:
            [\dots]}
            \par\hfill\textit{Kontext: \cref{sec:umsetzung-implementierung}}

    \item \enquote{Erkläre, welche Fehlerbilder verteilter Systeme der
            nachfolgende Ausführungspfad noch nicht abfängt, und schlage je
            Fall die konservativste Behandlung vor: [\dots]}
            \par\hfill\textit{Kontext: \cref{sec:umsetzung-implementierung}}
\end{enumerate}

\subsection{Öffentliches Dashboard (Cloud-Run-Anwendung)}
\label{sec:prompts-fable-dashboard}
\begin{enumerate}
    \item \enquote{Entwirf einen zustandslosen Dienst, der redigierte
            JSON-Feeds aus einem Objektspeicher liest und daraus eine
            Weboberfläche rendert, ohne dass die handelnden Maschinen einen
            eingehenden Port öffnen müssen: [\dots]}
            \par\hfill\textit{Kontext: \cref{sec:umsetzung-implementierung}}

    \item \enquote{Erzeuge die Oberfläche des Dashboards so, dass der
            Indexverlauf mit den Durchschnittslinien und den Regimewechseln
            dargestellt wird und je Entscheidung das aggregierte Votum mit
            aufklappbaren Einzelbegründungen erscheint: [\dots]}
            \par\hfill\textit{Kontext: \cref{sec:umsetzung-implementierung}}

    \item \enquote{Prüfe den folgenden Rendering-Pfad darauf, ob Inhalte aus
            dem Journal ungeprüft in das Dokument geschrieben werden, und
            schlage eine Ausgabekodierung vor: [\dots]}
            \par\hfill\textit{Kontext: \cref{sec:umsetzung-implementierung}}
\end{enumerate}

\subsection{Verifikation der erzeugten Testskripte}
\label{sec:prompts-fable-testverifikation}

Die Testskripte wurden mit Claude Opus~5 entworfen
(\cref{sec:prompts-claude-tests}); die nachfolgenden Prompts dienten allein
ihrer unabhängigen Gegenprüfung durch ein zweites Modell.

\begin{enumerate}
    \item \enquote{Prüfe die folgende Testsuite kritisch: Welche Verzweigung
            des Produktivcodes wird von keinem Testfall erreicht, und welcher
            Test würde auch dann bestehen, wenn die geprüfte Zusicherung
            entfiele? [\dots]}
            \par\hfill\textit{Kontext: \cref{sec:umsetzung-implementierung}}

    \item \enquote{Benenne Testfälle, die nur scheinbar prüfen, was ihr Name
            behauptet, und schlage jeweils eine Formulierung vor, die den
            tatsächlich geprüften Sachverhalt trifft: [\dots]}
            \par\hfill\textit{Kontext: Anhang~\ref{sec:anhang-archiv-tests}}
\end{enumerate}

\subsection{Erzeugung der KI-generierten Abbildungen}
\label{sec:prompts-fable-abbildungen}

Beide Abbildungen sind in der jeweiligen Bildunterschrift als KI-generiert
ausgewiesen (Leitfaden Kap.~5.7 in Verbindung mit Kap.~6.2.4). Inhaltliche
Grundlage waren ausschließlich der eigene Quellcodebestand und die eigene
Architekturplanung; die Modellantwort lieferte die grafische Umsetzung.

\begin{enumerate}
    \item \enquote{Erzeuge eine SVG-Grafik der Zielarchitektur des
            Handelsagenten auf Basis der bereits erfolgten Implementierungen
            sowie der geplanten Bestandteile; umgesetzte Elemente
            durchgezogen, geplante gestrichelt, mit Legende. Style die Grafik
            attributbasiert ohne \texttt{<style>}-Block, damit die
            Inkscape-Konvertierung verlustfrei arbeitet: [\dots]}
            \par\hfill\textit{Kontext: \cref{fig:zielarchitektur}}

    \item \enquote{Erzeuge aus der nachfolgenden ASCII-Repräsentation der
            Verzeichnis- und Dateistruktur eine SVG-Grafik, die die Struktur
            in einer monospaced Schriftart schwarz auf weiß rendert, die
            Ebenen-Header durch Fettsatz hervorhebt und je Datei die
            Herkunftsmaschine kennzeichnet: [\dots]}
            \par\hfill\textit{Kontext: \cref{fig:projektstruktur}}
\end{enumerate}

%  =============================================================================
%  Ende
%  =============================================================================

%  =============================================================================
%  Datei-Hinweise & Konfiguration: siehe config.md
%  -> Abschnitt "anhang/ki-prompts/tool-gemini.tex"
%  =============================================================================

\section{Google Gemini (Flash~3.5)}
\label{sec:prompts-gemini}

\kiwerkzeug{Google Gemini, Modell Flash 3.5}{Stand: Mai 2026}{https://gemini.google.com}
% Eintrag im KI-Verzeichnis (Leitfaden Kap. 6.2.4) erzwingen mittels \nocite.
\nocite{geminiFlash2026}

\subsection{Literaturrecherche}
\label{sec:prompts-gemini-literatur}
\begin{enumerate}
    \item \enquote{Nenne einschlägige, zitierfähige Quellen (Monografien,
            Normen, begutachtete Fachartikel) zur BPMN-2.0-Prozessmodellierung
            und gib jeweils vollständige bibliografische Angaben an: [\dots]}
            \par\hfill\textit{Kontext: \Cref{cha:grundlagen}}

    \item \enquote{Fasse den aktuellen Forschungsstand zur
            Geschäftsprozessmodellierung im ERP-Umfeld zusammen und benenne
            die zentralen Referenzwerke: [\dots]}
            \par\hfill\textit{Kontext: \Cref{cha:grundlagen}}
\end{enumerate}

\subsection{Erstellung des Abkürzungsverzeichnisses}
\label{sec:prompts-gemini-abkuerzungen}
\begin{enumerate}
    \item \enquote{Extrahiere alle Abkürzungen aus dem folgenden Text und
            erstelle ein alphabetisch sortiertes Abkürzungsverzeichnis mit
            jeweiliger Langform: [\dots]}
            \par\hfill\textit{Kontext: gesamte Arbeit}
    \item \enquote{Überprüfe das vorhandene Abkürzungsverzeichnis auf
            Vollständigkeit und Konsistenz und ergänze fehlende Einträge: [\dots]}
            \par\hfill\textit{Kontext: gesamte Arbeit}
\end{enumerate}

%  =============================================================================
%  Ende
%  =============================================================================

%  =============================================================================
%  Datei-Hinweise & Konfiguration: siehe config.md
%  -> Abschnitt "anhang/ki-prompts/tool-kimi.tex"
%
%  BESONDERHEIT DIESES ABSCHNITTS: Die Prompts wurden in ENGLISCHER Sprache
%  gestellt und sind hier in deutscher Uebersetzung wiedergegeben. Der Hinweis
%  steht sowohl im Vorspann des Abschnitts als auch an jedem einzelnen Prompt
%  ("aus dem Englischen uebersetzt") -- beides beibehalten, damit kein Prompt
%  fuer sich betrachtet als deutschsprachige Eingabe missverstanden wird.
%
%  ABSCHNITTSNUMMER: Dieser Abschnitt ist Anhang III.IV. Das \note-Feld des
%  Bib-Eintrags moonshotKimi2026 verweist darauf. Wird die Reihenfolge der
%  \input-Zeilen in anhang/ki-prompts.tex geaendert, ist der Verweis dort
%  mitzufuehren.
%  =============================================================================

\section{Moonshot AI Kimi (K3)}
\label{sec:prompts-kimi}

\kiwerkzeug{Moonshot AI Kimi, Modell K3}{Modellbezeichner
\texttt{kimi-k3}, Modellvorstellung vom 17.07.2026}{https://www.kimi.com}
% Eintrag im KI-Verzeichnis (Leitfaden Kap. 6.2.4) erzwingen mittels \nocite.
\nocite{moonshotKimi2026}

Eingesetzt für ein allgemeines Codereview der beigelegten Quellcodebeilage
(\cref{cha:anhang-archiv}) sowie für den Erstentwurf des
Änderungsverzeichnisses der Beilage. Das Modell erhielt ausschließlich das
entpackte Archiv und keinen Zugriff auf den Projektordner; es war damit die
einzige Prüfinstanz, die den Code in genau der Form vorfand, in der er einem
Dritten vorliegt.

Die Prompts dieses Abschnitts wurden in englischer Sprache gestellt und sind
nachfolgend sinngleich ins Deutsche übersetzt wiedergegeben; der Zusatz
\textit{aus dem Englischen übersetzt} weist an jedem Prompt darauf hin. Die
Übersetzung war nicht Teil der Eingabe an das Modell.

Wie bei den übrigen Werkzeugen wurde jede Antwort eigenständig gegengeprüft.
Das Ergebnis fiel hier geteilt aus und ist deshalb ausdrücklich vermerkt: Von
zwei vorgeschlagenen Änderungen wurde die erste nach eigener Messung
übernommen -- sie betraf den Selbstschutz des Sitzungserhalts und beschrieb
einen zuvor unentdeckten Fehler --, die zweite als Rückschritt verworfen. Ihr
lag eine Verwechslung zweier Spalten des Archivmanifests zugrunde, die ohne
Kenntnis des Bauprozesses nicht erkennbar war; ihre Übernahme hätte den Bau
der Beilage abgebrochen. Auch der Erstentwurf des Änderungsverzeichnisses
wurde vor der Übernahme in drei Punkten korrigiert.

\subsection{Allgemeines Codereview der Quellcodebeilage}
\label{sec:prompts-kimi-codereview}
\begin{enumerate}
    \item \enquote{Anbei der vollständige Quellcodebestand eines
            Handelsagenten in der Fassung, die einer wissenschaftlichen
            Arbeit als Beilage beigefügt wird. Prüfe ihn ohne Vorgabe eines
            Schwerpunkts auf Fehler und benenne je Befund die betroffene
            Datei, die Zeile, die Wirkung im Betrieb und deine Sicherheit über
            den Befund: [\dots]}
            \par\hfill\textit{Kontext: \cref{cha:anhang-archiv}; aus dem
            Englischen übersetzt}

    \item \enquote{Beschränke die Prüfung nun auf die Shell-Skripte und dort
            auf die Behandlung von Rückgabewerten: Welche Abfrage deutet einen
            Fehlerausgang als regulären Zustand, und in welchem Fall bleibt
            eine notwendige Handlung dadurch unbemerkt aus? [\dots]}
            \par\hfill\textit{Kontext: \cref{cha:anhang-archiv},
            \cref{sec:umsetzung-implementierung}; aus dem Englischen
            übersetzt}

    \item \enquote{Lege zu jedem übernahmewürdigen Befund eine minimale
            Änderung als Textunterschied vor, die die bestehende Struktur
            wahrt, und benenne für jede Änderung, in welchen Umgebungen sie
            das Verhalten ändert und in welchen sie nachweislich wirkungslos
            bleibt: [\dots]}
            \par\hfill\textit{Kontext: \cref{cha:anhang-archiv}; aus dem
            Englischen übersetzt}
\end{enumerate}

\subsection{Entwurf des Änderungsverzeichnisses der Beilage}
\label{sec:prompts-kimi-changelog}
\begin{enumerate}
    \item \enquote{Formuliere zu den bestätigten Befunden einen Abschnitt für
            den Anhang, der die Unterschiede zwischen der vorliegenden und der
            nächsten Version der Beilage benennt: je Punkt das fehlerhafte
            Verhalten davor, die Änderung und den Beleg, dass sie wirkt --
            sachlich, ohne Wertung und ohne Wiedergabe von Quelltext:
            [\dots]}
            \par\hfill\textit{Kontext:
            Anhang~\ref{sec:anhang-archiv-identitaet}; aus dem Englischen
            übersetzt}

    \item \enquote{Prüfe den Abschnitt gegen die Angaben der Beilage: Welche
            Zahl, welcher Dateiname und welcher Verweis darin sind aus dem
            beigefügten Bestand nicht belegbar? [\dots]}
            \par\hfill\textit{Kontext:
            Anhang~\ref{sec:anhang-archiv-identitaet}; aus dem Englischen
            übersetzt}
\end{enumerate}

%  =============================================================================
%  Ende
%  =============================================================================


%  =============================================================================
%  Ende
%  =============================================================================
